\chapter{Nominal Syntax and Reference}\label{ch:nominal-syntax}\label{ch:nouns}

\section{Overview}
\label{sec:nm-overview}

Iridian nouns distinguish direct, transitive, and indirect case in singular
and plural forms, with endings and stem grades distributed across nine
lexical declension classes. A restricted vocative survives in names, kinship
terms, and a small residue of other nouns. Case and number receive cumulative
exponence, while reference may remain unmarked, be anchored in the shared
discourse, or be presented as the speaker's selection of a particular referent.

Shared anchoring follows the case-number form or is realised cumulatively
with it; speaker anchoring has the proclitic exponent \ird{si}. Prenominal
material precedes the noun in the following order:

\begin{center}
(\ird{si}) + (\textsc{dem}) + (\textsc{quant}/\textsc{num}) +
(\textsc{poss}.\Ind{}) + (\textsc{prop.root} / \textsc{stative} / \ird{-ový}) +
\textsc{root} + (\textsc{case.num}) + (\textsc{anch})
\end{center}

The two anchoring positions are alternative realisations of the same nominal
category: \ird{si} and shared anchoring do not co-occur on one nominal spine,
although a possessor has its own anchoring independently of the containing
noun phrase (§\,\ref{sec:nm-si-complex}). The template records relative
positions rather than independently segmentable suffixes in every cell;
it also leaves the finer ordering of quantificational material unspecified.
The structural analysis of these positions is developed in
§\,\ref{sec:nm-funcseq}.

\subsection{The root and NP}
\label{sec:nm-root}

The lexical noun occupies the head position of \synhead{NP}. Unlike the verbal
root, which is licensed by a selecting functional head (the Asp.Align morpheme
is obligatory on every finite verb; cf.\ §\,\ref{sec:vm-root}), the nominal root
may surface without an overt case-number ending in consonant-stem and bare-root
classes, whereas thematic classes retain an overt direct ending. A nominal
without either \ird{si} or shared anchoring retains the extended structure
although \textsc{Anch\textsuperscript{0}} receives no overt spellout; its
referent may nevertheless be identifiable in context
(§\,\ref{sec:nm-bare-noun}).

Iridian argumental lexical nominals are analysed as projecting beyond
\synhead{NP}. The label \synhead{DP} serves as a cover term for the complete
extended nominal expression, while \synhead{AnchP} names its particular
D-layer and \synhead{NP} the lexical constituent below number, case, and
anchoring. A demonstrative or quantifier can therefore extend the expression
above AnchP without changing the use of DP as a category label. In this
terminology, a \emph{bare nominal} lacks overt anchoring; it need not be a
syntactically unexpanded NP. \emph{Definiteness} is the traditional
descriptive label for one common use of the anchoring layer.

The uniform extended structure is an analysis of Iridian nominal
distribution, including the unanchored lexical pivots of existential clauses
(§\,\ref{sec:existential-presentatives}); it is not inferred merely from
the possibility of a definite interpretation. DP status does not itself
require an individual-denoting expression or a silent uniqueness operator,
since the semantic composition of Anch is property-preserving
(§\,\ref{sec:nm-anchor-types}). This terminology leaves the minimum
structure selected within postpositional complements to the separate
analysis in §\,\ref{sec:nm-postp-structure}.

Absence of an overt article does not independently settle nominal category:
Despić develops an analysis without D for articleless nominals
\cite[§§\,1, 3]{despic2019}, whereas Borik and Espinal argue for DP
structure in Russian bare definite-kind expressions using additional
syntactic diagnostics \cite[§\,3.2]{borikespinal2019}. These alternatives
motivate keeping the structural analysis of Iridian distinct from the
interpretation of any single unmarked nominal.

\FloatBarrier

\section{Anchoring and reference}
\label{sec:nm-def}

A bare nominal leaves its referential interpretation to the construction and
context, whereas \ird{si} explicitly presents a particular referent selected
by the speaker. Suffixal or cumulative shared anchoring instead marks that
referent as accessible within the shared discourse. The contrast between
\ird{byl}, \ird{si byl}, and \ird{bylo} therefore concerns the basis on which
reference is presented, while the form of the exponent depends on the
feature bundle and, for shared anchoring, the noun's inflectional class.
The grammar represents these possibilities as alternative specifications of
Anch\textsuperscript{0}:

\medskip
\begin{center}
\begin{tabularx}{0.88\textwidth}{l l L}
  \toprule
  \textsc{feature bundle} & \textsc{exponent} & \textsc{description} \\
  \midrule
  {[+anch, +shr]} & suffixal or cumulative & recover the intended referent through the shared discourse model \\
  {[+anch, \textminus{}shr]} & \ird{si} & present the intended referent as the speaker's selection \\
  {[\textminus{}anch]} & $\varnothing$ & impose neither anchoring instruction; interpretation supplied by construction and context \\
  \bottomrule
\end{tabularx}
\end{center}
\medskip

These semantic instructions are independent of the number of segments that
realise them: a cumulative ending contributes the same shared anchoring as
a separate suffix. The analysis locates the contrast in nominal
interpretation rather than treating the three bundles as lexical
translations of English \trsl{the}, \trsl{a certain}, and \trsl{a}.

\subsection{Anchored reference and the [\textsc{+anch}] feature}
\label{sec:nm-def-semantics}

Shared-anchored nouns present a referent as recoverable within the common
discourse model, whether it functions as topic or remains within the comment.
Anaphoric uses resume a previously introduced referent, while situational
uses identify an entity salient in the speech situation; overt anchoring may
also reintroduce or single out a discourse-relevant token with a contrastive
or resumptive interpretation. Situational identifiability also permits an
unmarked nominal, so the availability of a uniquely identifiable referent
does not by itself require overt shared anchoring.

Suffixal and cumulative forms realise \textsc{Anch\textsuperscript{0}}[\textsc{+anch,+shr}]:
the shared feature instructs the addressee to recover the referent from the
common discourse model. Speaker-anchoring \ird{si}, which shares
[\textsc{+anch}] but lacks [\textsc{+shr}], instead presents the speaker's
selection of a particular referent (§\,\ref{sec:nm-spec}). The feature
[\textsc{+anch}] alone therefore does not require shared identifiability.

Shared accessibility concerns the recovery of the intended referent, which
may be established through prior mention or the speech situation; it does
not require that the noun's descriptive content have only one satisfier in
every relevant domain. A discourse containing several children can thus
provide the background against which a particular child is identified,
although the anchor alone supplies no means of choosing between equally
accessible candidates. Conversely, a contextually unique child can be
referred to by bare \ird{byl}, without the conventional retrieval
instruction contributed by \ird{bylo}.

Uniqueness and anaphoric familiarity are distinct sources of definite
interpretation in Schwarz's comparative account
\cite[§§\,2.1, 2.3--2.4]{schwarz2019}, while Coppock and Beaver distinguish
definiteness from individual-denoting determinacy and existential import
\cite[§§\,1--2]{coppockbeaver2015}. These distinctions supply analytical
diagnostics for Iridian without equating its anchoring bundles with the
articles or denotations of either account.

The grammar accordingly treats anchoring as a condition on the presentation
and identification of nominal reference, independently of the semantic
operation that makes the nominal an argument or a predicate. An overt anchor
does not by itself establish an individual-denoting semantic type, an
assertion that an individual exists, or a fixed scope relative to negation,
modality, or quantification. Where a construction asserts existence, as in
an existential presentation (§\,\ref{sec:existential-presentatives}), that
assertion belongs to the construction; any identification condition on its
nominal expression must be distinguished from the asserted content.

\subsection{Unmarked reference and distributional restrictions}
\label{sec:nm-bare-noun}
\label{sec:nm-def-restrictions}

A nominal without either \ird{si} or shared anchoring projects an
\synhead{AnchP} whose head lacks [\textsc{+anch}] and receives no overt
Vocabulary Item at PF. Its reference may be indefinite, contextually definite,
or generic; absence of overt anchoring leaves speaker identifiability
unmarked rather than requiring an unidentified or unindividuated referent.
Under this analysis, [\textsc{\textminus{}anch}] contributes no conventional
speaker-selection or shared-retrieval condition. It does not itself supply
existential closure, a uniqueness operator, or kind formation: the containing
construction and contextual interpretation determine which of these routes
is available. A silent Anch head is consequently compatible with a definite
reading without being identified with a silent definite article in every use.
The absence of overt anchoring thus corresponds to a feature distinction within the same
nominal structure, rather than to a smaller phrase. Where
[\textsc{+anch,+shr}] is present, Vocabulary Insertion supplies either a separate anchor or the
class-conditioned cumulative case-number--anchoring exponent (§\,\ref{sec:nm-def-nondirect}).

The vocative excludes overt shared anchoring (§\,\ref{sec:nm-vocative}).
Generic and kind-denoting reference ordinarily uses the bare noun, whose null
anchor leaves the referent's interpretation to the construction and context;
overtly anchored nominals generally resist generic readings. The ordinary
shared-anchored uses identify discourse tokens, while the eligibility of an
individuated kind for overt shared retrieval remains a separate descriptive
question (§\,\ref{sec:nm-kinds-mass}).

Shared anchoring has the suffixal and cumulative realisations described in
§\,\ref{sec:nm-def-allomorph}; their shared-accessibility interpretation
remains independent of the nominal expression's clausal role.

\subsection{Kinds, mass reference, and institutional interpretation}
\label{sec:nm-kinds-mass}

Bare \ird{vlak} \trsl{wolf} occurs as the kind-denoting topic of
\ird{Vlak drže na hlučna} \trsl{The wolf howls at night}
(\ref{ex:tc-generic-topic}). This established use combines nominal kind
reference with a generic predication; absence of an overt anchor alone does
not determine either component. An episodic reference to one wolf and a
reference to the wolf kind require different contextual interpretations even
when the nominal form remains bare. The preference for unmarked generic
nominals therefore establishes no universal incompatibility between kind
reference and a D-layer: Despić's and Borik and Espinal's treatments derive
bare kind interpretations under different structural assumptions
\cite{despic2019,borikespinal2019}. Whether an Iridian overt anchor can
retrieve a kind introduced as a discourse referent, as distinct from a
particular animal, is not established by the bare generic example.

Class~VIII \ird{pivo} \trsl{beer} provides an independent contrast between
mass and count interpretation. Its direct plural \ird{piva} selects
servings, varieties, or other countable instances
(§\,\ref{sec:nm-num}); shared-anchored direct singular \ird{pivó} and
plural \ird{pivá} preserve the morphological number distinction
(§\,\ref{sec:nm-def-allomorph}). The contrast between \ird{pivo} and \ird{pivó}
holds mass-noun morphology constant while varying shared anchoring; that
between \ird{piva} and \ird{pivá} does the same within the countable plural.
An English translation such as \trsl{the beer} does not determine whether
the intended referent is a contextually delimited quantity, a variety, or
some other discourse value. The contextual conditions licensing those
readings must be distinguished from the established availability of the
anchored forms. The atomic count-noun cardinal schema in
§\,\ref{sec:nm-anchor-types} consequently supplies no automatic analysis
of mass measurement or mass-kind reference.

Institutional activity interpretations require separate diagnostics from the
generic and situationally unique uses just described. English
\trsl{go to the hospital} can describe a conventional institutional activity
without requiring identification of one particular hospital; these
\emph{weak definites} are distinct from the \emph{weak-article definites}
discussed by Schwarz \cite[§§\,1--2, 4.1]{williams2019}. Their availability with
Iridian bare nominals is unestablished, and neither a null Anch head nor a
contextually unique building supplies that activity interpretation by itself.

\subsection{Discourse indices and situation-relative interpretation}
\label{sec:nm-anchor-indices}

Bare \ird{byl} and shared-anchored \ird{bylo} can identify a child already
available in the discourse, although only the latter conventionally requires
shared retrieval. The grammar represents this requirement by an index
\(i\), whose value \(g(i)\) is the intended referent under an assignment
\(g\). The index is analytical notation rather than an additional pronounced
morpheme. Prior mention can establish its value, while a salient participant
in the speech situation can also be represented in the common discourse
model without a preceding linguistic antecedent. The shared condition thus
concerns accessibility at the point of interpretation, rather than obligatory
second mention.

Let \(P_s\) be the nominal property evaluated relative to a resource
situation \(s\), and let \(c\) supply the discourse context. The analysis
represents shared anchoring as the following restriction on that property:
\[
 \mathrm{Anch}_{+\mathrm{shr},i}^{c,g}(P_s)
   = \lambda x[P_s(x)\land x=g(i)].
\]
This expression is defined only where \(g(i)\) is available and recoverable
through the shared discourse model in \(c\). It returns a restricted
property; an operation in the containing construction must still supply
argumental reference or existential closure where required. Successful
reference also requires the recovered value to satisfy the nominal
description. The accessibility condition does not assert that the restricted
property has a satisfier, and it cannot create a suitable referent merely
because the suffix is present. Situation-derived accessibility is compatible
with the established situational uses of shared anchoring; the conditions
under which a previously unavailable index can be accommodated remain
unspecified.

An unmarked nominal can instead receive a contextually definite reading
when its description identifies a unique entity in the relevant situation,
or resume an accessible discourse referent without conventionally imposing
the indexed retrieval condition. Situation restriction and indexed retrieval
can therefore converge on the same individual. Neither their overlap nor
the availability of a bare anaphoric reading eliminates the contrast in
conventional presentation. The situation variable restricts the descriptive
domain independently of the choice of an overt anchor, following the
distinction between situation-relative uniqueness and indexed anaphoricity
in Schwarz's analysis \cite[§§\,2.3--2.4]{schwarz2019}; Iridian does not
identify these mechanisms with mutually exclusive bare and suffixal forms.

Relational interpretation introduces a further dependency. In a description
of someone's head, the possessor must be resolved before the resulting
property of heads can identify a referent (§\,\ref{sec:nm-implicit-poss}).
A house--door or book--author bridging context likewise supplies a relation
to an antecedent rather than identity with that antecedent: retrieving a book
does not itself retrieve its author. Schwarz's German contrasts distinguish
part--whole from producer bridging \cite[§\,2.3, examples~(11)--(12)]{schwarz2019},
but no corresponding division between Iridian anchoring forms follows from
those contrasts. Whether a relationally inferred referent supports bare
reference, overt shared retrieval, or both depends on its interpretation in
Iridian; the index of a possessor and the index of the possessed referent
must remain distinct.

A fixed assignment \(g\) represents reference to one discourse value; it
does not by itself derive covariation with a quantifier. For an interpretation
schematically expressed by English \trsl{Every child who bought a book read
the book}, the relevant book must vary with the child--book dependency.
Such an interpretation requires assignments \(g_d\), resource situations
\(s_d\), or both to vary with the quantified cases \(d\), together with a
binding or discourse-update mechanism that makes the appropriate value
accessible. These are distinct analytical possibilities: situation binding
can vary a uniquely identifying description, while dynamic binding can vary
an anaphoric index \cite[§\,2.4 and note~10]{schwarz2019}. The English
schema identifies the dependency to be explained; it supplies no Iridian
acceptability judgement or evidence that shared anchoring licenses such
covariation. Nominal shared retrieval also remains independent of the
clause's aboutness relation (§\,\ref{sec:tc-anchoring-interface}).

\subsection{Semantic types and nominal composition}
\label{sec:nm-anchor-types}

The contrast between bare \ird{byl}, speaker-anchored \ird{si byl}, and
shared-anchored \ird{bylo} can be represented without requiring Anch to
saturate its nominal complement. In the analysis adopted here, each bundle
takes a property of type \(\langle e,t\rangle\) and returns a property of
the same type, giving Anch the type
\(\langle\langle e,t\rangle,\langle e,t\rangle\rangle\);
\(e\) includes plural individuals as well as atoms. At a
fixed evaluation situation \(s\), the three entries are:
\[
\begin{aligned}
 \mathrm{Anch}_{-\mathrm{anch}}(P_s)
   &= P_s,\\
 \mathrm{Anch}_{+\mathrm{shr},i}^{c,g}(P_s)
   &= \lambda x[P_s(x)\land x=g(i)],\\
 \mathrm{Anch}_{-\mathrm{shr},v}^{c}(P_s)
   &= \lambda x[P_s(x)\land x=v_{c,s}].
\end{aligned}
\]
The first entry is the identity function on properties. The second has the
shared-accessibility condition stated in §\,\ref{sec:nm-anchor-indices}.
The third is defined where \(v_{c,s}\) represents a witness selected by the
speaker relative to the relevant information state; it imposes no
shared-retrieval requirement. Its witness can be a group, as in established
plural \ird{si býle}. The parameter records speaker selection without
introducing existential quantification over choice functions. Neither it
nor the shared index guarantees that its value satisfies \(P_s\); a
mismatching value leaves an empty restricted property. The situation
subscript on \(v\) permits evaluation within an intensional dependency
without requiring selection to vary there. Whether the witness remains
constant or varies with that dependency is a scope and interpretation
question, which the entry and the PF position of \ird{si} do not settle.

A property supplied to an ordinary universal quantifier can therefore remain
unsaturated through Anch. Writing \(A\) for that property and \(R\) for
the predicate in its scope, the schematic quantifier has type
\(\langle\langle e,t\rangle,\langle\langle e,t\rangle,t\rangle\rangle\):
\[
 \mathrm{EVERY}(A)(R)=\forall x[A(x)\mathbin{\rightarrow}R(x)].
\]
This representation illustrates the composition of the established bare
quantified nominals in §\,\ref{sec:nm-quant}; it does not assign identical
distributivity or number semantics to \ird{nět}, \ird{niže}, and
\ird{nách}. Converting \(A\) to an individual before applying this
quantifier would create a type mismatch. Conversely, these Anch entries
cannot take an already generalised-quantifier-denoting \synhead{QP} of type
\(\langle\langle e,t\rangle,t\rangle\): the quantifier must compose with
the property returned by Anch, while the reverse composition would require
an independently justified type lift. Where the construction instead
requires individual reference, a subsequent definite shift can yield
\(\iota x A(x)\), defined when exactly one relevant individual satisfies
\(A\); existential composition with a scope predicate can yield
\(\exists x[A(x)\land R(x)]\). These operations occur after the anchoring
restriction, rather than being lexical contributions of the suffix or
\ird{si}; separating them follows the distinction between nominal
predication and argumental type shifts in Coppock and Beaver's account
\cite[§\,1]{coppockbeaver2015}.

Type compatibility alone does not license an overt anchor with every
quantifier. Restriction by \(x=g(i)\) or \(x=v_{c,s}\) admits at most one
individual, even when that individual is plural, so a universal applied
directly to this property would quantify over at most that one value.
It would not thereby range over all members of an independently salient set
of children. A set-based domain restriction, schematically
\(\lambda x[P_s(x)\land C(x)]\), has different truth conditions from
identity with an indexed individual. Etxeberria and Giannakidou analyse
such non-saturating restrictions in quantifier constructions
\cite[§§\,2.1, 3]{etxeberriagiannakidou2019}; their account supplies a
comparative alternative, rather than evidence for a second Iridian Anch
entry. The distribution and interpretation of overt anchoring with universal
quantifiers remain unestablished.

Numerals likewise differ from generalised quantifiers in the semantic
operation they supply. In \ird{hroná lobek} \trsl{three apples}, singular
noun morphology does not prevent the phrase from denoting a plurality.
A schematic cardinal modifier can form the property
\(\lambda X[{}^{*}P_s(X)\land |\operatorname{Atoms}(X)|=3]\), where
\({}^{*}P_s\) contains the nonempty sums of atomic apples and
\(\operatorname{Atoms}(X)\) is the set of atomic members of \(X\)
\cite[pp.\,15, 24--25]{marti2020}.
The numeral supplies the cardinality condition independently of the singular
Vocabulary Item (§\,\ref{sec:nm-quant}). In a putative speaker-anchored
cardinal phrase, selection of a three-apple plurality must consequently be
distinguished from selection of one apple before cardinal modification.
The current hierarchy and the lower singleton-restricting entry do not by
themselves derive selection of a particular three-apple group; no
acceptability or scope claim about \ird{si} with cardinals follows from the
linearisation template.

An intersective modifier combines with an unsaturated nominal property by
intersection, whereas a possessor first supplies the relevant relational
argument (§\,\ref{sec:nm-implicit-poss}). These operations leave a property
available for anchoring without requiring an individual to be selected
prematurely. The higher adnominal demonstrative can identify a deictically
located value satisfying the property returned by Anch; with shared anchoring,
that value must also coincide with the recoverable indexed referent.
This compatibility requirement preserves the independent contributions of
\ird{tak byl} and \ird{tak bylo} (§\,\ref{sec:nm-def-double}) without
deriving a general emphatic meaning from type composition. It supplies no
licensing condition for the unestablished combination of a demonstrative
with \ird{si}.

\subsection{Case and the ordering \synhead{AnchP} $>$ KP}
\label{sec:nm-def-stacking}

Suffixal anchoring follows the case-inflected noun, as in
\ird{byla}/\ird{bylato}. Where case-number and anchoring have separate
exponents, their morphological order is:

\begin{center}
\textsc{root} $+$ \textsc{case.num} $+$ \textsc{anch}
\end{center}

\noindent The order in \ird{byl-a-to} places the case-number exponent
inside shared anchoring. Under the compositional analysis developed in
§\,\ref{sec:nm-funcseq}, this supports \synhead{AnchP} above \synhead{KP}.
Anchoring is available across the structural cases and preserves the noun's
case-dependent clausal role (§\,\ref{sec:nm-case-system}), while its
referential contribution concerns the nominal as a whole. This distribution
motivates treating anchoring separately from case licensing, although it
does not uniquely establish the relative height of the two heads. The
grammar retains the following hierarchy, including in cells whose
case-number and shared-anchoring features receive cumulative exponence:

\begin{center}
\begin{forest}
  syntax tree
  [\synhead{AnchP}
    [\synhead{KP}
      [\synhead{NumP}
        [\synhead{NP}
          [{\textit{root}}]
        ]
        [{Num\textsuperscript{0}\\{\footnotesize number}}]
      ]
      [{K\textsuperscript{0}\\{\footnotesize case}}]
    ]
    [{Anch\textsuperscript{0}\\{\footnotesize anchoring}}]
  ]
\end{forest}
\end{center}

\noindent The tree represents the syntactic hierarchy before its PF
realisation, omitting the higher left-peripheral layers. Demonstratives are
distinct edge material outside \synhead{AnchP}.
Anch\textsuperscript{0}[\textsc{+anch,\,-shr}] receives left-edge
precedence within the extended nominal domain and is realised as \ird{si};
its position in the syntactic spine is unchanged
(§\,\ref{sec:nm-si-linearisation}). Suffixal and cumulative shared
anchoring are realised at the right edge of the nominal inflectional domain.

The overt suffixal order parallels the verbal domain, where lower verbal
morphology precedes the outer finite suffixes (§\,\ref{sec:vm-asp}).
Cumulative exponence reduces the number of independently segmentable endings
without changing the interpretation or hierarchical order of the nominal
heads.

\subsection{Speaker anchoring with the proclitic \ird{si}}
\label{sec:nm-spec}

In a speaker-anchored expression, the speaker has a particular
individual or group in mind without presenting its identity as already shared
with the addressee. The feature [\textsc{\textminus{}shr}] withholds a
conventional requirement to retrieve that identity through the shared
discourse; it does not assert that the addressee is incapable of identifying
the referent independently, or that the referent is absent from the common
ground. Speaker identifiability is an
established dimension of indefinite interpretation: Kagan's analysis of
Russian distinguishes it from free-choice effects and narrow scope relative
to a sentence-level operator \cite{kagan2011}. The Iridian analysis uses that
distinction without assigning \ird{si} a fixed quantifier scope.

Speaker selection concerns the intended witness, whether an individual or a
group, rather than the number of entities satisfying the nominal
description. Although an ordinary affirmative assertion about that witness
may commit the speaker to its existence, the feature bundle itself neither
asserts existence in the actual world nor fixes the nominal's interpretation
within an intensional or negative construction. The distributions of those
constructions must establish the available readings; the initial position of
\ird{si} supplies no independent scope diagnostic
(§\,\ref{sec:nm-si-linearisation}).

\pex
\a
\begingl
\gla byl//
\glb child.\Dir{}//
\glft \trsl{a child; the child; child} \lingcontext{interpretation supplied by context}//
\endgl
\a
\begingl
\gla si byl//
\glb \Sanch{} child.\Dir{}//
\glft \trsl{a certain child (that I have in mind)}//
\endgl
\a
\begingl
\gla bylo//
\glb child.\Dir{}.\Def{}//
\glft \trsl{the child} \lingcontext{explicit shared anchoring}//
\endgl
\xe

\noindent The bare noun can refer to a particular child known to the speaker,
to a child whose identity is unspecified, or to the kind in an appropriate
generic construction. Neither speaker knowledge nor ignorance is encoded by
the absence of an anchor. In \ird{si byl}, speaker selection is overt, but
\ird{si} remains optional even when the speaker can identify the child.
First-mention introductions therefore permit a useful distinction: a bare
nominal introduces a child without a conventional commitment about the
speaker's knowledge, while \ird{si byl} presents the speaker as having a
particular child in mind. Shared \ird{bylo} asks the addressee to recover the
child through the common discourse model.

Within a single nominal spine, the incompatibility of \ird{si} with shared
anchoring is represented by their alternative feature specifications on the
same Anch head. Vocabulary Insertion selects the speaker-anchored proclitic or
the suffixal or cumulative shared-anchoring exponent according to that
specification. The independent
demonstrative layer does not itself make every combination felicitous: where
a demonstrative identifies the same referent for the addressee, its
combination with speaker-specific \ird{si} has no established interpretation.

\subsection{Nominal-domain linearisation and prosodic attachment}
\label{sec:nm-si-linearisation}

\ird{Si} precedes the material of its nominal expression, whereas the
suffixal anchor follows the case-number form. Iridian specifies a language-specific,
feature-conditioned PF linearisation instruction which accounts for this
asymmetry: the speaker-anchored terminal
receives precedence at the edge of its nominal domain before Vocabulary
Insertion supplies \ird{si}. Its syntactic position and interpretation remain
those of Anch\textsuperscript{0}. The distinction between a clitic's
structurally defined positioning domain and its phonological host has an
independent precedent in Klavans's treatment of clitic dominance, precedence,
and liaison \cite[§§\,1.2--1.3]{klavans1985}; the particular domain and
instruction below constitute the analysis of Iridian.

The interpretation in §\,\ref{sec:nm-anchor-types} assigns properties and
identification conditions to these feature bundles without changing their
exponents or pronunciation domains. Discourse indices receive no separate
spellout, and subsequent semantic type shifts introduce no additional Anch
terminal for this linearisation rule to order.

For an anchoring terminal $a$, let $D_a$ be the designated anchoring
pronunciation domain of the lexical noun whose case-number layer is selected
by $a$. Its upper boundary is the highest of the designated nominal
projections on that noun's spine: \synhead{AnchP}, the quantificational
field, and \synhead{DemP}, where present. This boundary designates a PF
domain rather than identifying it with every projection which may contain
the nominal expression. A larger nominal in which this expression is a possessor, a governing
postpositional phrase, and a containing clause lie outside $D_a$. A possessor therefore
has its own domain and anchoring terminal, although its ordered terminal
sequence occurs as a dependent inside the host noun's domain. This definition uses the
nominal spine to identify ownership; it does not identify a domain by whichever
noun or clitic happens to occur first in the string.

The PF fusion of number with K precedes linearisation; suffixal attachment
to the nominal host is specified separately in §\,\ref{sec:nm-funcseq}. PF then orders
embedded nominal domains before the domain containing them, retaining each
completed dependent as an ordered block. Within $D_a$, the
default order preserves the established precedence of demonstratives,
quantificational material, possessors, and the head noun, with K and Anch
following their complements. If this default sequence is
$\alpha\mathbin{\frown}\langle a\rangle\mathbin{\frown}\beta$, and $a$
bears [\textsc{+anch,\,\textminus{}shr}], nominal-domain linearisation assigns:
\[
  \operatorname{Lin}(D_a)
    = \langle a\rangle\mathbin{\frown}\alpha\mathbin{\frown}\beta.
\]
Here $\frown$ is concatenation, while $\alpha$ and $\beta$ are ordered
sequences that may contain completed dependent domains. The instruction
reorders the single occurrence of the terminal belonging to this nominal
spine; every other precedence relation is retained. An anchoring terminal
inside a possessor remains within that possessor's block, and $a$ receives no
second pronunciation at its syntactic source.

Vocabulary Insertion follows this ordering step, realising $a$ as \ird{si}.
The corresponding [\textsc{+anch,+shr}] terminal is subject to ordinary
head-final linearisation and receives either a separate suffix or a
cumulative exponent with the adjacent case-number terminal (§\,\ref{sec:nm-def-nondirect}); null anchoring receives
no overt exponent. Thus the feature-conditioned instruction changes neither
the case-number features nor the local conditioning of nominal exponence.
In the following schemata, $N$--K stands for the root and fused case-number
terminal before Vocabulary Insertion, while $M$ contains the ordered
prenominal material, including any completed possessor. The final strings
abbreviate the phonological realisation of that material:
\[
\begin{aligned}
  [_{D_a}\;M\;N\text{--K}\;a_{[+\mathrm{anch},-\mathrm{shr}]}]
    &\longmapsto a\;M\;N\text{--K}
      \longmapsto \textit{si}\;M\;N\text{--K},\\
  [_{D_a}\;M\;N\text{--K}\;a_{[+\mathrm{anch},+\mathrm{shr}]}]
    &\longmapsto M\;N\text{--K}\;a
      \longmapsto M\;N\text{--K--}\textit{anch}.
\end{aligned}
\]
The first arrow represents linearisation and the second Vocabulary Insertion;
\textit{anch} abbreviates shared anchoring; the second schema displays
the concatenative realisation, while specified cells realise K and Anch
cumulatively (§\,\ref{sec:nm-def-allomorph}). The minimal
speaker-anchored derivations yield the established \ird{si byl} and
\ird{si býle}, while their shared-anchoring counterparts yield \ird{bylo}
and \ird{býlete}. The variable $M$ specifies the placement instruction for
otherwise licensed nominal material, without establishing new combinations
of \ird{si} with demonstratives or quantifiers.

The instruction belongs to the mapping from nominal structure to PF
precedence, rather than to syntactic head movement. It therefore leaves the
hierarchy interpreted at LF, including the source of speaker anchoring,
unchanged. It also supplements the ordinary linearisation rules rather than
applying Local Dislocation across an intervening string: in the Distributed
Morphology account of Embick and Noyer, Local Dislocation is restricted to
linear adjacency \cite{embicknoyer2007}. No succession of adjacent swaps or
additional nominal head is required by the Iridian rule. The domain boundary
and the precedence instruction are explicit commitments of this analysis;
their availability does not follow independently from the general
architecture of Distributed Morphology.

Prosodification applies to the resulting string. \ird{Si} is an
unstressed proclitic outside the minimal lexical prosodic word
(§\,\ref{sec:stress}); its host is the following independently stressed
word. In \ird{si byl} and \ird{si býle}, the noun retains its initial lexical
stress; \ird{si} contributes no syllable to the host's subsidiary footing
(§\,\ref{sec:secondary-stress}). Where prenominal material intervenes,
the first independent word in that material supplies the host, which may belong to a possessor even though
that possessor has its own anchoring domain. The dependent's order is protected
during linearisation, but its initial word remains available for prosodic
attachment. Structural ownership, edge precedence, and phonological attachment
therefore contribute separately to the pronunciation of the nominal phrase.

\subsection{Complex nominals and nested anchoring domains}
\label{sec:nm-si-complex}

An attributive modifier intervenes between the speaker anchor and the noun
in the constructed phrase below, which combines \ird{si byl} with the
established modifier \ird{hradná} (§\,\ref{sec:nm-prop}):

\pex
\a
\begingl
\gla si hradná byl//
\glb \Sanch{} good.\textsc{attr} child.\Dir{}//
\glft \trsl{a certain good child (that I have in mind)}//
\endgl
\xe

\noindent The clitic precedes the modifier, whose initial lexical stress
is preserved; the noun is consequently neither its immediate host nor the
leftmost independent word of $D_a$.

A prenominal possessor creates two potential sources for the same initial
\ird{si}. The following constructed nominal schemata use the established
indirect \ird{bylam} and direct \ird{tóm}; their brackets identify anchoring
ownership rather than pauses or additional pronounced material:

\pex
\a
\begingl
\gla {[si bylam]} tóm//
\glb {\Sanch{} child.\Ind{}} book.\Dir{}//
\glft \trsl{a book belonging to a certain child I have in mind}
\lingcontext{speaker anchoring on the possessor}//
\endgl
\a
\begingl
\gla si {[bylam] tóm}//
\glb \Sanch{} {child.\Ind{} book.\Dir{}}//
\glft \trsl{a certain book I have in mind, belonging to a child}
\lingcontext{speaker anchoring on the containing nominal}//
\endgl
\xe

\noindent Both structures yield \ird{si bylam tóm} under nominal-domain
linearisation. In the first, the possessor's own terminal receives initial
precedence inside its dependent block; in the second, the containing
nominal's terminal precedes that block. The book's null anchor in the first
structure leaves its reference to context, just as the child's null anchor
does in the second. The resulting ambiguity is a prediction of the adopted
PF rule, rather than an independently established ambiguity judgement.
Although \ird{bylam} supplies the first independent prosodic host in both
structures, the location of stress does not identify the source of anchoring.

If both domains bear speaker anchoring, the same rule preserves both terminals
and predicts \ird{si si bylam tóm}. The feature alternatives on a single Anch
head do not exclude this structure, since the terminals belong to different
nominal spines. The rule prescribes no haplology or special clitic-cluster repair.

A speaker-anchored nominal may establish a new topic, occur as a focused or
backgrounded argument, or introduce an individual in an existential
presentation. Its subsequent omission depends on recoverability in the
continuing discourse, rather than on the presence of \ird{si} at first
mention (§§\,\ref{sec:tc-ordinary-topics}, \ref{sec:tc-topicless},
\ref{sec:tc-topic-drop}). The contrast between speaker selection and
aboutness is thus preserved when the same nominal expression receives
different clausal interpretations.

\FloatBarrier

\section{Demonstratives and quantification}
\label{sec:nm-upper}

\subsection{The demonstrative system and noun-internal DemP}
\label{sec:nm-dem}

Iridian demonstratives form a three-way deictic system with proximal, medial,
and distal values. The same paradigm has three distinct morphosyntactic
realisations: an adnominal demonstrative, a pronominal demonstrative, and an
indirect-case possessive demonstrative. The noun-internal analysis developed
here concerns the adnominal realisation, namely \synhead{DemP}, while the pronominal and possessive uses draw on the same
case-inflecting paradigm. Their positions within the noun phrase are set out
in §\,\ref{sec:nm-np-structure}.

\begin{table}[ht]
  \footnotesize
  \caption{Demonstrative paradigm relevant to noun-internal and
  pronominal syntax.}\label{tab:nm-dem-paradigm}
  \begin{tabularx}{0.7\textwidth}{L l l l}
    \toprule
                 & {\sc direct} & {\sc transitive} & {\sc indirect} \\
    \midrule
    Proximal     & \ird{tak}  & \ird{tam}  & \ird{ten} \\
    Medial       & \ird{tuto} & \ird{toho} & \ird{taly} \\
    Distal       & \ird{je}   & \ird{jaky} & \ird{jení} \\
    \bottomrule
  \end{tabularx}
\end{table}

\noindent Within the noun phrase, the adnominal demonstrative realises
\synhead{DemP}, the outermost noun-internal demonstrative projection above
\synhead{AnchP}. The hierarchy adopted for these distributions is
\synhead{DemP} $>$ \synhead{AnchP} $>$ \synhead{KP} $>$ \synhead{NumP} $>$ \synhead{PossP} $>$
\synhead{NP}, matching the descriptive order \textsc{demonstrative} $>$
\textsc{poss}.\Ind{} $>$ \textsc{head} given in §\,\ref{sec:nm-np-structure}.
Demonstratives stand outside the possessive grouping, whereas possessors remain
inside the \synhead{KP}/\synhead{NP} domain. Co-occurrence, as in \ird{tuto že tom}, shows directly that
demonstratives and possessors cannot be analysed as occupants of one and the
same structural slot.

\begin{center}
\begin{forest}
  syntax tree
  [\synhead{DemP}
    [{\textit{adnominal demonstrative}}]
    [\synhead{AnchP}
      [\synhead{KP}
        [\synhead{NumP}
          [\synhead{PossP}
            [{\textit{oblique possessor}.\Ind{} / $\varnothing$}]
            [{\synhead{NP}\\\textit{head noun}}]
          ]
          [{Num\textsuperscript{0}\\{\footnotesize number}}]
        ]
        [{K\textsuperscript{0}\\{\footnotesize case}}]
      ]
      [{Anch\textsuperscript{0}\\{\footnotesize anchoring features}}]
    ]
  ]
\end{forest}
\end{center}

The core semantics of the demonstrative layer is deictic: proximal indexes a
referent near the speaker, medial a referent near the addressee, and distal a
referent away from both. The topic-proximal, addressee-linked medial, and
obviative or backgrounded distal uses (§\,\ref{sec:nm-dem-pron}) are
pragmatic extensions of those lexical deictic values.
The formal analysis therefore keeps lexical deixis and discourse tendency
distinct.

The indirect demonstratives
\ird{ten}, \ird{taly}, and \ird{jení} may occupy the ordinary possessor slot
before the noun, as in \ird{ten dum}. In such forms the demonstrative is not the
outer deictic determiner of the noun phrase; it is an indirect-case possessor
inside \synhead{PossP}. The grammar therefore allows two distinct noun-internal
positions fed by the same demonstrative system: adnominal \synhead{DemP} above
\synhead{AnchP}, and demonstrative possessors inside \synhead{PossP}.

\subsection{Double-marking: demonstrative and shared anchoring}
\label{sec:nm-def-double}

A demonstrative may co-occur with overt shared anchoring, as in
\ird{tak bylo} \trsl{this very child; this particular child here}
(§\,\ref{sec:nm-def}). Whereas the demonstrative contributes deixis, the
shared anchor contributes discourse identifiability. Their combination
permits emphatic, resumptive, contrastive, or strongly individuating readings;
none of these readings follows obligatorily from the two markers alone.

The theoretical analysis adopted here is that demonstratives occupy a
\synhead{DemP} above \synhead{AnchP} (§\,\ref{sec:nm-dem}). Co-occurrence is
grammatical because the two layers are distinct; it is marked because the noun
phrase is being both deictically located and additionally singled out as the
anchored discourse token. Multiple marking does not independently establish
semantic redundancy: the domain-restricting analyses discussed by Etxeberria
and Giannakidou illustrate how superficially repeated definiteness marking
can make distinct contributions \cite[§§\,4.1--4.2]{etxeberriagiannakidou2019}.
For Iridian, the additional indexed-retrieval condition is separated from the
choice among these contextual readings (§\,\ref{sec:nm-anchor-types}).

\pex
\a
\begingl
\gla tak byl//
\glb \Dem{}.\Prox{} child.\Dir{}//
\glft \trsl{this child}//
\endgl
\a
\begingl
\gla tak bylo//
\glb \Dem{}.\Prox{} child.\Dir{}.\Def{}//
\glft \trsl{this very child; this particular child here}//
\endgl
\xe

\subsection{Pronominal demonstratives and third-person reference}
\label{sec:nm-dem-pron}

Iridian lacks independent third-person personal pronouns. The full demonstrative
paradigm therefore serves as the ordinary third-person referential series
(§\,\ref{sec:nm-dem}). This pronominal use is not identical to noun-internal
\synhead{DemP}, because no lexical head noun is present inside the expression,
but it is derived from the same underlying demonstrative system and must be kept
distinct from the noun-internal ordering facts summarised in
§\,\ref{sec:nm-np-structure}.

In discourse, proximal forms commonly track the current topic or most accessible referent, medial forms
a secondary or addressee-linked referent, and distal forms a backgrounded,
contrastive, or obviative referent. In formal terms, these effects are best
treated as interpretive extensions available to pronominal demonstratives, not
as separate lexical paradigms (§\,\ref{sec:nm-dem}).

\ex
\begingl
\glpreamble \lingcontext{Constructed sequence contrasting three discourse referents.}//
\gla Tak ščének, tuto pjášček, je zemácek.//
\glb \Prox{}.\Dir{} arrive-\Pf{}, \Med{}.\Dir{} eat-\Antip{}-\Pf{}, \Dist{}.\Dir{} run-\Pf{}//
\glft \trsl{He arrived, he ate, that one ran.}//
\endgl
\xe

\subsection{The reduced anaphoric series}
\label{sec:nm-dem-reduced}

Alongside the full paradigm, Iridian has a reduced anaphoric series
\ird{i}/\ird{ji}/\ird{ni}, corresponding to distal
\ird{je}/\ird{jaky}/\ird{jení}. These forms are restricted to pronominal use,
only when the referent is already highly accessible in the discourse, and they
are excluded as ordinary adnominal demonstratives of the type described in
§\,\ref{sec:nm-np-structure}.

Semantically and pragmatically, the reduced series signals weak anaphoric
resumption rather than full deictic pointing. Formally, it belongs to the
demonstrative system, but it is not the spellout of noun-internal
\synhead{DemP}. A natural DM treatment is to analyse \ird{i}/\ird{ji}/\ird{ni}
as specialised pronominal Vocabulary Items competing with the full distal series
under a discourse-anaphoric high-accessibility condition; see again
§\,\ref{sec:nm-dem}.

\pex
\a \ird{je} \trsl{that one; fully deictic distal}
\a \ird{i} \trsl{he/she; weak anaphoric resumption}
\xe

\subsection{Numerals and the quantificational field}
\label{sec:nm-quant}

Iridian places numerals and quantifiers prenominally. They belong to the upper
nominal field, but the presently available evidence does not force a
fine-grained cartographic split among cardinals, adnominal interrogatives, and
universal quantifiers or negative indefinites. The grammar treats them as
occupants of a single quantificational modifier field between the demonstrative
domain and the lower \synhead{NP}-internal material.

Cardinal numerals require the noun within their own nominal phrase to appear
in its \emph{singular} form, regardless of the quantity expressed. Singular nominal morphology under
cardinals has parallels in Hungarian, Finnish, and Turkish
\cite[§\,1]{marti2020}, although this comparison concerns number marking
rather than the case constructions of those languages. In Iridian, the cardinal
specifies the quantity denoted by the phrase, while the noun's case-number slot
is realised by the singular Vocabulary Item.

\pex
\a \ird{hroná lobek} \trsl{three apples} \hfill\textit{(cardinal $+$ \textsc{sg})}
\a \ird{hroná lúbce} \hfill\textit{(cardinal $+$ \textsc{pl}, ungrammatical)}
\xe

\noindent In DM terms, the fused [K$+$Num] head inside a cardinally-quantified
\synhead{DP} is spelled out by the singular Vocabulary Item: the Q head
(\synhead{QP}) saturates the number semantics of the expression, so
Num\textsuperscript{0} within \synhead{KP} receives no [+pl] feature to trigger
a plural VI. This morphological specification leaves the cardinal phrase's
plural interpretation intact; semantic number composition is distinguished
from argumental saturation in §\,\ref{sec:nm-anchor-types}.

The exact-cardinality interpretation is supplied by the numeral, while
approximate readings depend on intonation and contextual hedging. The
singular requirement applies within the numeral's own nominal phrase;
it does not extend into an independent partitive complement
(§\,\ref{sec:nm-ext-indirect}). In \ird{ona pivum štava}
\trsl{one mug of beer}, the numeral quantifies the mug, while indirect
\ird{pivum} specifies its contents.

Adnominal interrogatives pattern with the same prenominal field. \ird{Jak}
selects from a contextually restricted set, whereas the other interrogatives
simply introduce the relevant open variable. In determiner use they modify the
noun phrase directly rather than occupying the clause-level \emph{wh}-position
themselves.

\pex
\a \ird{jede ruščevní svirkošc} \trsl{what Russian author}
\a \ird{jak ruščevní svirkošc} \trsl{which Russian author}
\xe

Universal quantifiers and negative indefinites can likewise be adnominal: \ird{nět študent}
\trsl{all students}, \ird{niže jec} \trsl{all dogs}, \ird{nách študent}
\trsl{each student}, \ird{žé tóm} \trsl{not one book}. Negative indefinites,
including pronominal \ird{deno} \trsl{nothing}, require clausal negative
licensing (§\,\ref{sec:nm-negative-concord}); their adnominal use restricts
the existential variable to the following noun's denotation. Negative
existence is expressed by the existential predicate
(§\,\ref{sec:neh}), whose suppletive \ird{zad-} realises the same clausal
polarity dependency.

\FloatBarrier

\section{Pronouns, anaphors, and nominal operators}
\label{sec:nm-pron}

\subsection{Personal pronouns as full DPs}
\label{sec:nm-pron-personal}

First- and second-person pronouns are best treated as fully projected DPs whose
nominal content is supplied by person features rather than by a lexical noun
root. They inflect for the same three productive cases as lexical nouns; their
indirect forms serve as oblique possessors in the analytic possessive
construction described in §\,\ref{sec:nm-poss}. Iridian lacks a separate
third-person personal-pronoun paradigm; that function is taken over by the
demonstrative system described in
§§\,\ref{sec:nm-dem}--\ref{sec:nm-dem-reduced}.

\begin{table}[ht]
  \footnotesize
  \caption{Personal pronouns.}
  \label{tab:nm-personal-pronouns}
  \begin{tabularx}{0.65\textwidth}{L l l l}
    \toprule
                  & {\sc direct} & {\sc transitive} & {\sc indirect} \\
    \midrule
    1\textsc{sg}  & \ird{dá}   & \ird{dam}  & \ird{že}   \\
    2\textsc{sg}  & \ird{já}   & \ird{jí}   & \ird{jena} \\
    1\textsc{pl}  & \ird{mé}   & \ird{mně}  & \ird{mi}   \\
    2\textsc{pl}  & \ird{tová} & \ird{tvy}  & \ird{tě}   \\
    \bottomrule
  \end{tabularx}
\end{table}

The indirect case forms serve as the possessors in analytic possessive
expressions: possession is structurally encoded through the indirect
domain, and the possessor occupies \synhead{PossP} as a prenominal indirect-case
\synhead{DP} (§\,\ref{sec:nm-poss}).

\pex
\a \ird{že dum} \trsl{my house}
\a \ird{mi dum} \trsl{our house}
\xe

\subsection{The reflexive as local nominal anaphor}
\label{sec:nm-pron-refl}

The reflexive is not a discourse-tracking pronoun but a true anaphor. Its
interpretation depends on a local antecedent in the clausal domain, and it
competes with demonstrative third-person forms precisely by forcing coreference
where demonstratives would permit disjoint reference.

\begin{table}[ht]
  \footnotesize
  \caption{Reflexive anaphor forms.}
  \label{tab:nm-reflexive}
  \begin{tabularx}{0.5\textwidth}{L l}
    \toprule
    {\sc case} & {\sc form} \\
    \midrule
    Direct     & \ird{svem} \\
    Transitive & \ird{svě}  \\
    Indirect   & \ird{svy}  \\
    \bottomrule
  \end{tabularx}
\end{table}

\pex\label{ex:nm-reflexive-connectivity}
\a
\begingl
\gla Svem Marku védnik.//
\glb \Refl{}.\Dir{} Marek.\Trs{} see-\textsc{trs}-\Pf{}//
\glft \trsl{Marek saw himself.}//
\endgl
\a
\begingl
\gla I Marku védnik.//
\glb 3.\Dist{}.\Dir{} Marek.\Trs{} see-\textsc{trs}-\Pf{}//
\glft \trsl{Marek saw him.}//
\endgl
\xe

\noindent The reflexive reading in (\ref{ex:nm-reflexive-connectivity}a)
can be indexed as \ird{svem}\textsubscript{$i$}
\ird{Marku}\textsubscript{$i$}: the actor is the antecedent of the
direct-case patient. In (\ref{ex:nm-reflexive-connectivity}b),
\ird{i}\textsubscript{$j$} instead permits a distinct referent from
\ird{Marku}\textsubscript{$i$}. These indices specify the readings under
discussion; they are not nominal exponents or a claim that every demonstrative
must have disjoint reference.

The actor's thematic position is structurally above the patient complement
of V (§\,\ref{ch:verbs}). The local relation needed for the reflexive can
therefore be represented independently of the pronounced patient--actor order:
\[
 [_{\mathrm{VoiceP}}\ \mathrm{actor}_{i}
   [_{\mathrm{VP}}\ \mathrm{svem}_{i}\ V]\ \mathrm{Voice}].
\]
This is a representation of argument structure, not an additional surface
order. The reflexive requires connectivity with its local antecedent even
when it precedes that antecedent. Such connectivity does not, by itself,
establish obligatory reconstruction in every interpretive respect. Ordinary
placement pronounces the patient above the actor while preserving the lower
occurrence at which its local anaphoric dependency is interpreted
(§\,\ref{sec:th-particle-framework}); the pronounced order creates no new
antecedent relation. The distribution
across causative, passive and nominal domains requires construction-specific
antecedent conditions; discourse prominence alone does not supply them.

The intensifying \trsl{own} reading of a reflexive is analysed as an extension
of local coreference within the anaphoric domain, without a separate possessive
adjective.

\subsection{Interrogatives and quantificational operators}
\label{sec:nm-pron-operators}

Interrogatives, indefinites, and universal quantifiers can stand pronominally
or modify a noun phrase adnominally, while the negative series also includes
locative and temporal expressions. Adnominal forms pattern with the
quantificational field discussed in §\,\ref{sec:nm-quant}; in pronominal use they project as DPs without an overt lexical noun.

\begin{table}[ht]
  \footnotesize
  \caption{Core interrogative nominal operators.}
  \label{tab:nm-interrogatives}
  \begin{tabularx}{0.8\textwidth}{L L L L}
    \toprule
    \ird{jede} & who & \ird{jak} & which \\
    \ird{ježe} & what & \ird{zajehu} & why \\
    \ird{jehát} & whom & \ird{jiká} & how many \\
    \ird{jehu} & how & \ird{jišká} & how much \\
    \ird{jemí} & when & \ird{jeně} & to where \\
    \ird{jena} & where & \ird{jení} & from where \\
    \bottomrule
  \end{tabularx}
\end{table}

\begin{table}[ht]
  \footnotesize
  \caption{Selected negative indefinites and corresponding universal expressions.}
  \label{tab:nm-indefinites}
  \begin{tabularx}{0.8\textwidth}{L L L L}
    \toprule
    {\sc negative form} & {\sc meaning} & {\sc universal form} &
    {\sc meaning} \\
    \midrule
    \ird{zide} & no one & \ird{nět} & everyone \\
    \ird{deno} & nothing & \ird{niže} & everything \\
    \ird{žena} & nowhere & \ird{nina} & everywhere \\
    \ird{žemě} & never & \ird{nimě} & always \\
    \ird{žé} & not one & \ird{nách} & each \\
    \bottomrule
  \end{tabularx}
\end{table}

The negative series consists of dependent
indefinites whose interpretation includes a nominal, locative, temporal, or
cardinality restriction. The universal series retains universal or
distributive quantificational force and does not bear the negative-concord
feature described below.

Wh-hood ---
the property of introducing an open variable --- is a lexical or
\synhead{DP}-internal feature of the interrogatives: they are operators
that may occupy argument or adjunct positions inside the clause. Questionhood
proper, by contrast, is clausal and is located in \synhead{PolP}: the clause is
typed as a question by the Pol\textsuperscript{0} head (overt \ird{by} or null),
not by the presence of a wh-operator in Spec,\synhead{FocP} (see
§\,\ref{sec:by-pol-head} in Chapter~\ref{ch:clause-types} and
§\,\ref{sec:th-wh-interrogation} in Chapter~\ref{ch:clause-types}). A wh-expression participates in ordinary argument or adjunct placement
without dedicated wh-fronting, or moves to Spec,\synhead{FocP} for marked
focus at the edge above the ordinary placement field. Its lower chain
occurrence supplies the interrogative variable within the domain of Pol;
neither ordinary placement nor marked focus alone determines whether the
clause is interrogative. The same
interrogative nominals listed in Table~\ref{tab:nm-interrogatives} appear both
in matrix content questions and in embedded non-assertive \ird{poby}-complements
(§\,\ref{sec:poby} of Chapter~\ref{ch:embedding}), where they function as
wh-operators inside an issue-denoting embedded clause rather than as
direct-question operators.

\pex
\a \ird{jede bych ščenžkou?} \trsl{who arrived yesterday?}
\a \ird{jak ruščevní svirkošc} \trsl{which Russian author}
\a \ird{nět študent} \trsl{all students}
\a \ird{zide byl} \trsl{no child} \lingcontext{adnominal phrase; requires clausal negative licensing}
\xe

The negative series requires local clausal negative licensing; its distribution,
examples and analysis are given in §\,\ref{sec:nm-negative-concord}.

\FloatBarrier

\section{Noun-phrase structure and modification}
\label{sec:nm-np-structure}

\subsection{The functional sequence and linear order}
\label{sec:nm-funcseq}

The canonical linear order of the Iridian noun phrase is:

\begin{center}
\ird{si} $>$ \textsc{dem} $>$ \textsc{quant/num} $>$
\textsc{poss}.\Ind{} $>$ \textsc{prop.root / stative / -ový} $>$
\textsc{head} + \textsc{case.num} + \textsc{anch}
\end{center}

The sequence combines a head-final suffixal spine with prenominal
modifiers and a speaker anchor whose position is determined by nominal-domain
linearisation (§\,\ref{sec:nm-si-linearisation}). It records relative slots
without establishing the finer ordering or interpretation of every possible
combination of quantifiers and demonstratives.

\pex
\a \ird{tak študencí tóm} \trsl{this student's book}
\a \ird{jak ruščevní svirkošc} \trsl{which Russian author}
\xe

The nominal expression is analysed as an extended functional projection
under the \synhead{DP} Hypothesis (Abney 1987). Anch\textsuperscript{0}
expresses shared anchoring through suffixal or cumulative exponents and
speaker anchoring through \ird{si}; the traditional label \emph{definiteness}
therefore identifies a descriptive part of a wider referential system.
Table~\ref{tab:nm-projections} sets out the projections from the lexical
root to the demonstrative domain.

\begin{table}[ht]
  \footnotesize
  \caption{Functional projections of the Iridian \synhead{DP} (innermost to
  outermost).}\label{tab:nm-projections}
  \begin{tabularx}{\textwidth}{l L L l}
    \toprule
    {\sc position} & {\sc morpheme/item} & {\sc syntactic projection} &
    {\sc see} \\
    \midrule
    Root          & lexical stem        & \synhead{NP} / $\surd$P
                  & §\,\ref{sec:nm-root} \\
    Property root & small closed class  & \synhead{PropP} (adjunct to \synhead{NP})
                  & §\,\ref{sec:nm-prop} \\
    Possessor     & oblique noun phrase & \synhead{PossP} (prenominal possessor layer inside \synhead{KP})
                  & §\,\ref{sec:nm-poss} \\
    Case $\times$ Number & six fused case-number endings & \synhead{KP} $>$ \synhead{NumP};
                       K\textsuperscript{0} and Num\textsuperscript{0} fuse at PF
                  & §\,\ref{sec:nm-case}, §\,\ref{sec:nm-num} \\
    Anchoring     & suffixal or cumulative
                       (Anch\textsuperscript{0}[\textsc{+anch,+shr}]); \ird{si}
                       (Anch\textsuperscript{0}[\textsc{+anch,\,-shr}],
                       proclitic)
                  & \synhead{AnchP} (D-layer; shared anchoring suffixal or cumulative,
                       \ird{si} left-edge proclitic)
                  & §\,\ref{sec:nm-def}, §\,\ref{sec:nm-spec} \\
    Quantification & numerals, adnominal interrogatives,
                     universal quantifiers and negative indefinites
                   & quantificational modifier field in the upper
                     nominal domain; exact cartographic slot left
                     conservative
                   & §\,\ref{sec:nm-quant}, §\,\ref{sec:nm-spec} \\
    Demonstrative & proximal \ird{tak}/\ird{tam}/\ird{ten};
                       medial \ird{tuto}/\ird{toho}/\ird{taly};
                       distal \ird{je}/\ird{jaky}/\ird{jení}
                  & \synhead{DemP} (noun-internal realisation of the
                       demonstrative system)
                  & §\,\ref{sec:nm-dem} \\
    \bottomrule
  \end{tabularx}
\end{table}

The sequence \textsc{root--case.num--anch} places case-number inside
shared anchoring, consistently with the hierarchy \synhead{AnchP} $>$
\synhead{KP} $>$ \synhead{NumP}. The working interpretation of the
Mirror Principle in §\,\ref{sec:vm-overview} relates morphological
constituency to syntactic composition, subject to independently specified
realisation mechanisms \cite{baker1985,harley2011}. The suffixal sequence
therefore supports this hierarchy under the assumed compositional mapping;
it cannot uniquely determine the hierarchy independently of that mapping.
The referential distribution discussed in §\,\ref{sec:nm-def-stacking}
provides a separate reason for distinguishing anchoring from case licensing.

Head-final linearisation orders the complement of each inflectional head
before that head, whereas formation of a morphological word requires an
additional attachment relation; affixation and head movement are themselves
distinct properties \cite{harley2013b}. For Iridian, PF realisation combines
the lexical nominal host with the immediately following [K$+$Num] exponent
and, where present, the separate shared-anchoring exponent into one word.
This attachment applies only to the contiguous nominal inflectional sequence;
it neither incorporates prenominal dependents nor crosses into a possessor's
nominal domain. The schematic morphological constituency is
$[[N\text{--}\mathrm{K.Num}]\text{--}\mathrm{Anch}]$, although class-selected
cumulative exponents may realise more than one of these structural positions
(§\,\ref{sec:nm-def-allomorph}). Suffixal attachment is thus an explicit
realisation assumption, rather than a consequence of head-finality alone.

Num\textsuperscript{0} and K\textsuperscript{0} are distinct heads in the
syntax and fuse at PF before Vocabulary Insertion (§\,\ref{sec:nm-num});
zero exponence and conditioned stem grades do not remove either projection.
The higher prenominal material has its own distribution, with
\synhead{DemP} realised at the left edge of its domain. Speaker anchoring
likewise retains its source in Anch\textsuperscript{0}, while the rule in
§\,\ref{sec:nm-si-linearisation} assigns \ird{si} left-edge precedence and
proclitic attachment. Since this rule specifies a different mapping from
structure to pronunciation, the position of \ird{si} cannot be used as an
ordinary suffix-order diagnostic.

The suffixal nominal spine parallels the verbal complex: case-number precedes
anchoring within the noun phrase, just as the lower verbal morphology precedes
the outer finite suffixes. The following tree represents this nominal hierarchy before PF linearisation.

\begin{center}
\begin{forest}
  syntax tree
  [\synhead{DemP}
    [{\textit{demonstrative series}}]
    [\synhead{QP}
      [{\textit{numeral / quantifier}}]
      [\synhead{AnchP}
        [\synhead{KP}
          [\synhead{NumP}
            [\synhead{PossP}
              [{\textit{oblique possessor}}]
              [{\synhead{NP}\\\textit{root}}]
            ]
            [{Num\textsuperscript{0}\\{\footnotesize number}}]
          ]
          [{K\textsuperscript{0}\\{\footnotesize case}}]
        ]
        [{Anch\textsuperscript{0}\\{\footnotesize anchoring features}}]
      ]
    ]
  ]
\end{forest}
\end{center}

\noindent The tree represents the common nominal hierarchy.
Overt exponents of higher nominal projections are shown directly at the edge of
the projection they realise, and silent heads are not inserted merely to force a
specifier analysis. Num, K, and Anch are shown as separate syntactic heads
before PF realisation; the proclitic exponent of Anch
receives its precedence at linearisation rather than at the head-final
position shown here (§\,\ref{sec:nm-si-linearisation}).
The \synhead{QP} layer is drawn only as a conservative cover term for prenominal
quantificational material; the analysis does not assume a finer
numeral-vs-quantifier split unless the data require it.
Its placement above \synhead{AnchP} is compatible with property-taking
composition, but establishes neither the felicity of every quantifier--anchor
combination nor a set-based domain-restricting meaning for overt anchoring
(§\,\ref{sec:nm-anchor-types}).

\subsection{The oblique possessor in the DP}
\label{sec:nm-poss}
\label{sec:nm-poss-struct}

Iridian expresses possession by means of an oblique (indirect-case) noun phrase
in prenominal position (§\,\ref{sec:nm-ext-indirect}). The oblique possessor
\ird{Marcí} in \ird{Marcí dum} \trsl{Marek's house} occupies a structural
position immediately above the head \synhead{NP} and inside the \synhead{KP}
domain, below \synhead{AnchP} and below noun-internal \synhead{DemP}. In terms
of the \synhead{DP} hierarchy, this is the prenominal dependent position of
\synhead{PossP}.

The possessor bears the general indirect case, whose nominal uses include
material, source, and partitive dependence. The same morphological value also
marks core arguments in verbal clauses, including imperfective patients and
causative causees; indirect case therefore does not diagnose peripheral
argument status. Nominal relational licensing and clausal argument licensing
are distinct dependencies which receive the same K realisation.

\begin{center}
\begin{forest}
  syntax tree
  [\synhead{AnchP}
    [\synhead{KP}
      [\synhead{NumP}
        [\synhead{PossP}
          [{\textit{oblique possessor}.\Ind{}}]
          [{\synhead{NP}\\\textit{head noun}}]
        ]
        [{Num\textsuperscript{0}\\{\footnotesize number}}]
      ]
      [{K\textsuperscript{0}\\{\footnotesize case ending}}]
    ]
    [{Anch\textsuperscript{0}\\{\footnotesize anchoring features}}]
  ]
\end{forest}
\end{center}

\noindent The possessor appears before the head noun in the surface string,
consistent with head-final base-generation of specifiers. The entire possessive
phrase [\synhead{PossP} oblique-N head-N] forms a constituent that is then
inflected for number and case by Num and K and optionally for anchoring by
\textsc{Anch\textsuperscript{0}}.

\subsection{Implicit possessors and relational nominal interpretation}
\label{sec:nm-implicit-poss}

Body-part, kinship and part--whole nouns permit an unexpressed semantic
possessor when the relation is contextually recoverable. In a relational
interpretation, \ird{čás} \trsl{head} denotes a head of some individual:
\[
 \lambda p.\lambda h[\mathrm{HEAD}(h)\land\mathrm{PART\_OF}(h,p)].
\]
Licensing the implicit relation and identifying its referent are distinct.
The relational nominal permits the dependency to remain unexpressed;
independently licensed binding, accessible discourse referents and utterance
perspective constrain its resolution. This lexical-semantic distinction is
independent of nominal declension class.

\pex\label{ex:nm-implicit-poss}
\a
\begingl
\gla Čás konca.//
\glb head.\Dir{} hurt-\Cont{}//
\glft \trsl{My / his / her / your head hurts.}
\lingcontext{possessor resolved from the utterance context}//
\endgl
\a
\begingl
\gla Že čás konca.//
\glb 1\textsc{sg}.\Ind{} head.\Dir{} hurt-\Cont{}//
\glft \trsl{My head hurts.} \lingcontext{explicit possessor; clarification or contrast}//
\endgl
\a
\begingl
\gla Marcí čás konca.//
\glb Marek.\Ind{} head.\Dir{} hurt-\Cont{}//
\glft \trsl{Marek's head hurts.} \lingcontext{explicit nominal possessor}//
\endgl
\xe

In a spontaneous self-report, the first example favours the speaker's head.
A report about a patient or a narrative centred on Marek can instead favour
that individual. These are defeasible contextual readings: neither the noun
nor \ird{konca} carries first-person possessor agreement. The head is the
sole clausal argument in these examples. An overt possessor remains an
ordinary prenominal indirect-case dependent inside \synhead{PossP}; it is
available for clarification, contrast or a less accessible referent. It adds
no clausal experiencer, external possessor or applied participant.

Resolving the possessor first supplies a property of heads belonging to
\(p\); identifying a head satisfying that property is a further operation.
Thus \ird{že čás} and \ird{Marcí čás} fix the possessor overtly while
leaving the head noun's anchoring unmarked. A unique head relative to a
resolved possessor need not have been mentioned previously, and the
speaker-oriented interpretation of bare \ird{čás} neither supplies covert
\ird{si} nor establishes a shared index for that head.

The semantic argument \(p\) is not itself a claim that the nominal
contains a silent pronoun. Provisional \(pro_{\mathrm{POSS}}\) in the
existing \synhead{PossP}, inside \synhead{KP}, can represent the dependency;
direct contextual saturation of the relational argument remains a distinct
analysis. Neither representation entails mechanical deletion of overt
\ird{že}, and omission is not defined as first-person pronoun deletion.

An overt possessor or independently licensed binding dependency fixes the
relation before discourse preferences apply. A continuing aboutness topic
can favour an accessible possessor, but neither the most recent subject nor
the grammatical pivot determines that referent automatically. Several equally
salient candidates can make the bare nominal ambiguous and favour an overt
possessor without making omission categorically ungrammatical.
Converbial same-pivot continuity likewise imposes no possessor-identity rule
(§\,\ref{sec:switch-reference}).

Ordinary alienables such as \ird{dum} \trsl{house} and \ird{tóm}
\trsl{book} need no obligatory lexical possessor argument. Their association
with an owner can be recovered where the context supplies it; a bare house
or book does not automatically mean the speaker's possession. Speaker
anchoring \ird{si} selects a referent, and shared anchoring marks shared
identifiability; neither specifies ownership or the person of an implicit
possessor (§\,\ref{sec:nm-def}). The possessed noun remains overt throughout:
this nominal dependency is separate from clausal argument omission and its
construction-specific licensing (§\,\ref{sec:ref-management}).

\subsection{Number within the possessive DP}
\label{sec:nm-poss-ne}
\label{sec:nm-poss-number}

A possessor and its containing nominal each bear their own fused case-number
form. The indirect plural \ird{býlac} therefore marks several children as
possessors without requiring the possessed object to be plural, while the
direct plural \ird{lúbce} marks several apples independently of the number
of their possessor (§\,\ref{sec:nm-num}). The following constructed pair
combines these established forms with indirect singular \ird{bylam} and
direct singular \ird{lobek}:

\pex
\a
\begingl
\gla býlac lobek//
\glb child.\Ind{}.\textsc{pl} apple.\Dir{}.\textsc{sg}//
\glft \trsl{an apple belonging to children}//
\endgl
\a
\begingl
\gla bylam lúbce//
\glb child.\Ind{}.\textsc{sg} apple.\Dir{}.\textsc{pl}//
\glft \trsl{apples belonging to a child}//
\endgl
\xe

\noindent Number is local to each nominal spine, just as anchoring is:
plurality in one does not require plurality in the other. Neither expression
contains an overt anchor, so contextual identifiability remains independent
of the number contrast. Speaker anchoring may likewise belong to the possessor
or to its containing nominal (§\,\ref{sec:nm-si-complex}).

\subsection{Stative verbs and property roots as \synhead{NP} modifiers}
\label{sec:nm-prop}

Ordinary property modification uses a dedicated non-finite \synhead{AttrP}
before the noun. ATTR selects a property projection or reduced state
predication, not a finite verb with its aspect ending removed. The lexical
property roots that also feed stative verbs belong to this lower property
domain. Prenominal property modification follows the possessive layer and
precedes the \synhead{NP}.

The head-final principle correctly predicts that all prenominal modifiers —
whether stative-verb-based or property-root-based — appear before the head noun,
never after it. Postnominal modification is not a productive construction in
standard Iridian; the prenominal restriction on modifiers is a direct
consequence of the head-final parameter.

\pex
\a \ird{múžna nuž} \trsl{a beautiful city}
\a \ird{hradná byl} \trsl{a good child}
\a \ird{svorná dum} \trsl{a big house}
\a \ird{všihňa byl} \trsl{a happy child}
\a \ird{žbirňa stol} \trsl{a wet table}
\xe

The listed root-strong Class~I \ird{múžna}, suffix-strong Class~II
\ird{hradná}/\ird{svorná}, and soft Class~III
\ird{všihňa}/\ird{žbirňa} realise this same ordinary attributive
function (§\,\ref{sec:stative-attr}). Class membership supplies neither
nominal definiteness nor agreement or temporal anchoring; it does not create
a universal contrast between enduring and episodic properties.

\noindent These ordinary attributives are restricted to adnominal use within
the \synhead{DP}. Clause-level manner modification is expressed instead by the
derived adverbs in \ird{-ě} discussed in §\,\ref{sec:stative-adv} of
Chapter~\ref{ch:derivational}; the attributive itself is not extended to ordinary
adverbial use.

\subsection{The \ird{-ový} attributive and its embedded DP}
\label{sec:nm-ovni}

The relative head \ird{-ový} selects the same voice--aspect
configuration as a participant nominal in \ird{-ou}, but forms a prenominal
modifier of an overt noun (§\,\ref{sec:th-relativization}). Its
\synhead{RelP} contains a silent operator linked to the privileged participant
gap in the embedded predicate; abstraction over that position supplies the
property intersected with the head noun's denotation. The modifier precedes the
noun, and its suffix does not agree with the noun in case.

The embedded configuration retains viewpoint and participant orientation,
while nominal relational heads license indirect-case thematic dependents
independently of matrix case. Thus perfective \ird{pjáščikový} describes an
agent and \ird{védnikový} a patient. The head noun receives case from its
matrix clause or, when it is itself a dependent within another relative,
from that containing nominal domain. It does not receive the case that its
silent counterpart would bear in a finite clause.

\pex
\a \ird{lúbcí pjáščikový pres} \trsl{the man who ate the apple}
\a \ird{[Jancí védnikový lúbcí] pjáščikový pres}
   \trsl{the man who ate the apple that Janek saw}
\xe

The operator in each relative corresponds to the participant privileged by its
own voice--aspect configuration. In the nested example, \ird{Jancí} is an
agent dependent of patient-oriented \ird{védnikový}, whereas \ird{lúbcí},
modified by that relative, is the theme dependent of agent-oriented
\ird{pjáščikový}. Neither relative assigns matrix case to \ird{pres}.
The exact attachment of \synhead{RelP} relative to possessive material remains
undetermined; its prenominal position and independent case domain do not depend
on that choice. An ordinary continuous property modifier uses the dedicated
attributive with quantity-conditioned \ird{-á}/\ird{-a}
(§\,\ref{sec:stative-attr}). A continuous stative relative in \ird{-ový}
is separately licensed by an explicit episodic, stage-level, temporal or
contrastive context: \ird{mužnový nuž} identifies a city beautiful at the
relevant stage, whereas \ird{múžna nuž} simply attributes beauty. The former
retains Asp and abstracts over the sole state holder; the latter is a reduced
property modifier without Asp. The
selected \synhead{AspP} includes inner valency, causative and orientation
structure, but ends below clausal Neg, Phas, the Comm-field, Mood and clausal
evidentiality (§\,\ref{sec:nominalised}). Thus the bracket in
\ird{[lúbcí pjáščikový] pres} contains the positive participant description;
placing \ird{vy}, \ird{so}, \ird{och}, \ird{kamo} or clausal negation
inside that reduced description is excluded. An operator in the matrix clause
may take scope over a predication involving the entire DP; it does not become
part of the relative. The restriction on reduced \ird{-ový} does not ban
negative finite relatives.

\subsection{The \ird{ty}-supplement}
\label{sec:nm-appos}

Iridian possesses an appositional particle \ird{ty} that introduces a
right-peripheral supplement: a phrase that is coreferent with a preceding
nominal and presents it as a clarification, identification, or elaboration. The
semantics are those of identity---\textit{that is}, \textit{namely}---and the
\ird{ty}-phrase is always postnominal, set off from its antecedent by a prosodic
break.

The \ird{ty}-phrase is not a restrictive modifier. It does not narrow the
denotation of the head noun the way an adjective or relative clause does: a
phrase like \ird{Karel ty mlazka} does not pick out a subset of Karels who are
brothers, it simply adds the information that the referent already identified as
Karel is the speaker's brother. The \ird{ty}-phrase is therefore best analysed
as a \term{supplement}---a nominally-anchored, non-restrictive phrase that is
coreferent with its antecedent but structurally peripheral to it, not integrated
into the \synhead{NP}'s extended projection. This distinguishes \ird{ty}-phrases
categorically from adjectives and relative modifiers, which are adjuncts within
the \synhead{NP} and compositionally restrict the head. Because \ird{ty}-phrases
are supplements rather than adjuncts, they are not governed by the head-final
principle: modifiers fall within the \synhead{NP} and must precede their head;
supplements are peripheral and their position---always to the right of the
antecedent---is determined by information structure and discourse organisation.

The antecedent nominal is introduced first; \ird{ty} then opens the supplement
phrase. When the supplement consists of a bare nominal, no further element is
required: the nominal is itself the head and the identity relation is immediate.

\pex
\a \ird{Karel ty mlazka}
   \trsl{Karel---that is, my brother}
\xe

When the supplement consists of non-nominal material---an adjective, a
\ird{-ový}-relative, a postpositional phrase---a pronominal head is required to
host it. Modifiers in Iridian must precede a nominal head; without one, the
modifier cannot form a well-formed phrase. This head is supplied by the
anaphoric pronoun \ird{i}, which appears at the right edge of the supplement and
corefers with the antecedent nominal.

\pex
\a \ird{nuž ty múžna i}
   \trsl{the city---that is, the beautiful one}
\a \ird{pres ty že védnikový i}
   \trsl{the man---that is, the one I saw}
\xe

\noindent The internal structure of the \ird{ty}-phrase respects head-finality
throughout: \ird{múžna} and \ird{že védnikový} precede their head \ird{i} in
the ordinary left-to-right modifier-before-head order. Head-finality is not
suspended in the supplement; it applies within it. The distribution of \ird{i}
follows directly: it is required whenever the \ird{ty}-phrase contains a
non-nominal modifier that needs a nominal head to attach to, and absent when the
\ird{ty}-phrase already contains a nominal head.

Ordinary attributive modification---adjectives, \ird{-ový}-relatives, and
similar prenominal material---remains the unmarked strategy for adnominal
description and carries no special discourse implication. The
\ird{ty}-supplement is marked for deliberate elaborative intent: it presents the
clarification as a distinct informational act, separate from the introduction of
the referent. In spoken language this implicature is weak, since production
pressures naturally favour establishing the referent before specifying it; in
formal and written language the implicature is stronger, and a \ird{ty}-phrase
reads as emphatic, contrastive, or rhetorically foregrounded. The construction
is not excluded from formal writing but is marked there, and is most natural for
heavy adnominal material---long nominalised relative clauses, for
instance---that would be unwieldy in prenominal position.

The full discourse-interactional properties of the \ird{ty}-supplement,
including its use as a repair and afterthought device in spoken interaction, are
treated in Chapter~\ref{ch:spoken} (\S\,\ref{sec:spoken-ty}).

\FloatBarrier

\section{Indirect nominal dependence and postpositions}
\label{sec:nm-ext}

\subsection{Material, partitive, source, and other \synhead{NP}-external
  dependencies}
\label{sec:nm-ext-indirect}

The indirect case is the general dependency form of the nominal
system. Possessors are simply its most tightly integrated \synhead{DP}-internal
use. The same K-form appears with material nouns, partitives, source
expressions, and other semantically dependent nominals that are not the
privileged head of the phrase. In this wider domain the indirect does
not mark a special inventory of separate morphological cases; rather, it
is the single oblique realisation of K whose interpretation is narrowed
by the lexical relation between nouns or by the head that selects the
dependent nominal.

\pex
\a \ird{kuní prosc} \trsl{silver spoon}
\a \ird{ona pivum štava} \trsl{one mug of beer}
\a \ird{Študencím na hroná} \trsl{three of the students}
\a \ird{Marcí z houba} \trsl{a gift from Marek}
\xe

In the partitive \ird{študencím na hroná}, the indirect plural
\ird{študencím} denotes the containing set of students, while \ird{na}
closes its postpositional phrase before the cardinal. The cardinal therefore
does not impose singular morphology on this separate nominal complement
(§\,\ref{sec:nm-quant}).

Material, partitive, and source dependents therefore stand on a
continuum with possessors. The possessor is the most tightly integrated
prenominal use of indirect \synhead{KP} inside the \synhead{DP}. Material and partitive
dependents are looser nominal relations. Postpositional phrases are the
next step outward: the same indirect nominal domain serves as the
complement of a relational head that closes the phrase from the right.

\subsection{Postpositional complements and phrase structure}
\label{sec:nm-postpositions}
\label{sec:nm-postp-structure}
\label{sec:nm-postp-attachment}

Iridian adpositions are head-final and are best analysed as
\term{postpositions}. The constant generalisation is that the nominal
complement bears indirect case. In the smallest representation, the
postposition combines with an indirect \synhead{KP}. In fuller noun phrases, that
\synhead{KP} may itself be embedded inside \synhead{AnchP} or higher nominal layers, but the
selectional requirement remains the same: the nominal core is oblique,
and the postposition closes the complement phrase at its right edge.

The indirect nominal complement combines with a following postposition whose
lexical meaning specifies a spatial, participant, source, privative, or
substitutive relation.

The structural asymmetry is straightforward. The nominal complement is
assembled first: the lexical/possessive domain combines with Num and
receives indirect case in K, and may then be expanded by anchoring or other higher nominal
material. Only after that nominal phrase is complete does the
postposition merge as the head of the larger postpositional phrase.
Head-finality then yields the surface order \textit{complement} +
\textit{postposition}.

\begin{center}
\begin{forest}
  syntax tree
  [PostP
    [\synhead{AnchP}
      [\synhead{KP}
        [\synhead{NumP}
          [\synhead{PossP}
            [{\textit{oblique possessor} / $\varnothing$}]
            [{\synhead{NP}\\\textit{head noun}}]
          ]
          [{Num\textsuperscript{0}\\{\footnotesize number}}]
        ]
        [{K\textsuperscript{0}\\\Ind{}}]
      ]
      [{Anch\textsuperscript{0}\\{\footnotesize anchoring features}}]
    ]
    [{Post\textsuperscript{0}\\\ird{na} / \ird{de} / \ird{še} / \dots}]
  ]
\end{forest}
\end{center}

\noindent The tree is intentionally conservative. What matters is not whether
every postpositional complement is always projected up to \synhead{AnchP}, but
that the complement is at least an indirect \synhead{KP} and that any higher
nominal material remains inside the complement phrase rather than intervening
between the complement and the postposition.

The postposition attaches to the right edge of the governed phrase as a whole,
not to the noun stem in isolation. A bare noun, an anchored noun, a possessed
noun phrase, and a noun phrase with prenominal modifiers all behave alike in
this respect: the nominal material comes first, and the postposition closes the
entire phrase.

\pex
\a \ird{dumě na} \trsl{in the house; at home}
\a \ird{bylamo my} \trsl{for the child}
\a \ird{Marcí dumě na} \trsl{at Marek's house}
\a \ird{svorná mostě nad} \trsl{over the big bridge}
\xe

\noindent The same logic extends to larger complements. Demonstratives,
possessors, property-root modifiers, stative modifiers, and \ird{-ový} phrases
all belong to the nominal complement if they modify the noun phrase itself. The
postposition does not split them apart or attach to the lexical noun before they
are assembled. This is the chief diagnostic for analysing the forms as
postpositions over phrasal complements rather than as suffixes on nouns or as
free adverbs loosely associated with them.

\subsection{Spatial and participant relations}
\label{sec:nm-postp-classes}
\label{sec:nm-postp-extended}

The common postpositions form a compact but semantically varied system. Their
lexical meanings differ, but their syntax is uniform.

\begin{table}[ht]
  \footnotesize
  \caption{Core postpositions and their main semantic classes.}
  \label{tab:nm-postpositions}
  \begin{tabularx}{\textwidth}{l l L}
    \toprule
    {\sc class} & {\sc members} & {\sc core relation} \\
    \midrule
    Locative & \ird{na} & location, interiority, surface contact \\
    Goal & \ird{de} & direction, endpoint \\
    Proximity & \ird{u} & nearness \\
    Accompaniment & \ird{še} & accompaniment, comitative association \\
    Benefactive & \ird{my} & benefit, intended recipient, purpose-oriented relation \\
    Source / separation & \ird{z} & source, origin, movement away, extraction \\
    Privative & \ird{nam} & absence \\
    Replacement & \ird{pale} & substitution, `instead' \\
    Exception & \ird{ahte} & exclusion, `except' \\
    \bottomrule
  \end{tabularx}
\end{table}
\FloatBarrier

\noindent The core spatial and participant relations are illustrated
below:

\pex
\a \ird{dumě na} \trsl{in the house}
\a \ird{vruham de} \trsl{to the town}
\a \ird{vruham u} \trsl{near the town}
\a \ird{bylam my} \trsl{for the child}
\a \ird{Jancí še} \trsl{with Janek}
\a \ird{Marcí z} \trsl{from Marek}
\a \ird{záhárí nam} \trsl{without sugar}
\a \ird{Jancí pale} \trsl{instead of Janek}
\a \ird{Jancí ahte} \trsl{except for Janek}
\xe

Benefactive \ird{bylam my} contains an indirect-case complement of a
postposition, whereas an applied beneficiary is a DP licensed within the verbal
argument structure. In a simple indicative applicative, the child takes direct
case, as in \ird{Byl Marcí travě pjaštébime} \trsl{Marek is eating bread
for the child's benefit} (§\,\ref{sec:vm-appl}). Both constructions can
express benefaction, although only the applicative selects that beneficiary
for the participant nominal and corresponding relative. A beneficiary retained
in indirect case under higher causative or directive orientation remains an
applied argument; its lack of \ird{my} and its licensing by Appl distinguish
it from the postpositional construction.

Alongside the broad locative and goal markers, Iridian also admits a set of more
specialised spatial and path postpositions. These behave like the core set in
selecting the same indirect nominal complement and in appearing at the right
edge of the phrase.

\pex
\a \ird{stolě pod} \trsl{under the table}
\a \ird{mostě nad} \trsl{over the bridge}
\a \ird{domě měz} \trsl{between the houses; among the houses}
\a \ird{vruham do} \trsl{up to the town}
\a \ird{leše mně} \trsl{through the forest; via the forest}
\xe

The same head-final syntax thus extends across broad location, bounded path,
vertical relation, interstitial relation, and traversal.

\subsection{Clausal distribution and prosodic attachment}
\label{sec:nm-postp-clausal}
\label{sec:nm-postp-status}

The postpositional phrase is not confined to nominal modification. The same
structure appears as a clausal adjunct, as the location complement of
existential predication, and as a dependent nominal expression inside larger
noun phrases. Its nominal complement retains the same internal structure in each use.

\pex
\a
\begingl
\gla Marek dumě na jeca.//
\glb Marek.\Dir{} house.\Ind{} in \Exst{}-\Cont{}//
\glft \trsl{Marek is at home.}//
\endgl
\a
\begingl
\gla Kroumaště na pivo so jeca.//
\glb refrigerator.\Ind{} in beer.\Dir{} still \Exst{}-\Cont{}//
\glft \trsl{There is still some beer left in the refrigerator.}//
\endgl
\a
\begingl
\gla Bylam my Jancí še stojáž.//
\glb child-\Ind{} for Janek-\Ind{} with go-\Ctpl{}//
\glft \trsl{(I am) going for the child with Janek.}//
\endgl
\a
\begingl
\gla Marcí z houba.//
\glb Marek-\Ind{} from gift.\Dir{}//
\glft \trsl{a gift from Marek}//
\endgl
\xe

\noindent The existential examples connect directly to §\,\ref{sec:jec} in
Chapter~\ref{ch:predicates}: the entity whose existence is asserted stands in the
direct case, while the locative PP contributes the place relation. More broadly,
postpositional phrases behave like ordinary head-final dependent phrases in the
clause. They precede the final predicate and preserve their internal noun-first,
head-last organisation.

The common postpositions attach prosodically to the preceding independent word
at the right edge of their nominal complement. They remain outside that word's
minimal lexical stress domain and ordinarily bear no primary stress of their
own, and their syllables are excluded from its subsidiary footing
(§§\,\ref{sec:stress} and~\ref{sec:secondary-stress}). In \ird{dumě na},
\ird{na} leans on \ird{dumě}; in \ird{Marcí dumě na}, its immediate host is again
\ird{dumě}, while \ird{Marcí} retains its separate lexical stress. The
syntactic complement is the whole nominal phrase, whereas the immediate
phonological host is its final independent word.

These \term{postposed clitic adpositions} are written separately. Prosodic
attachment neither requires orthographic fusion nor turns the postposition into
a nominal suffix. The construction combines a phrase that bears indirect case
with a head-final relational marker; its weak pronunciation and its phrasal
syntax describe different properties of the same element.
Head-final order and obligatory adjacency to the nominal complement distinguish
this construction from a free clause adverb.

\subsection{Postpositions and lexical preverbs}
\label{sec:nm-postp-preverbs}

The bound lexical preverbs \ird{de-} and \ird{mně-} share directional and
traversal meanings with postpositional \ird{de} and \ird{mně}, but attach
to a verbal root rather than selecting an indirect nominal complement
(§\,\ref{sec:lexical-preverbs}). Replacement and privation are expressed
by the preverbs \ird{pře-} and \ird{ne-} in verbal derivation, while the
corresponding postpositions remain \ird{pale} and \ird{nam}. These
synchronic relations do not establish a direction of historical derivation.
Accumulating \ird{na-}, as in \ird{navad-} \trsl{synthesise considerations
into a judgement}, likewise belongs to a verbal stem, whereas locative
\ird{na} follows an indirect nominal complement. Their formal resemblance
permits a hypothesis of common origin, but neither that hypothesis nor their
spatial associations identify the prefix with a nominal relational head.
Prefixal accumulation also supplies no comitative participant of the kind
expressed by \ird{še}; any applicative argument depends on the ordinary
verbal derivation rather than on \ird{na-} itself.

\FloatBarrier
